%%%%%%%%%%%%%%%%%%%%%%%%%%%%%%%%%%%%%%%%%%%%%%%%%%%%%%%
%%                                                   %%
%%    Стилевой файл для оформления курсовый работ    %%
%%                                                   %%
%%    Автор И.К. Марчевский                          %%
%%                                                   %%
%%%%%%%%%%%%%%%%%%%%%%%%%%%%%%%%%%%%%%%%%%%%%%%%%%%%%%%

\NeedsTeXFormat{LaTeX2e}
\ProvidesPackage{fn2kursstyle}[19/03/2022 V.2.3. fn2kursstyle. Marchevsky Ilia]

\newif\ifoglav\oglavfalse
\newif\ifspisok\spisokfalse
\newif\ifleftheadline\leftheadlinefalse
\newif\ifboldsect\boldsectfalse
\newif\ifeqwhole\eqwholefalse
\newif\iffigwhole\figwholefalse
\newif\iffootwhole\footwholefalse
\newif\ifkursrab\kursrabfalse
\newif\ifremarks\remarksfalse
\newif\ifhyperprint\hyperprintfalse
\newif\ifgreekit\greekitfalse
\newif\ifhyperref\hyperreffalse
\newif\ifnodotheadline\nodotheadlinefalse

\DeclareOption{oglav}{\oglavtrue}
\DeclareOption{spisok}{\spisoktrue}
\DeclareOption{leftheadline}{\leftheadlinetrue}
\DeclareOption{boldsect}{\boldsecttrue}
\DeclareOption{eqwhole}{\eqwholetrue}
\DeclareOption{figwhole}{\figwholetrue}
\DeclareOption{footwhole}{\footwholetrue}
\DeclareOption{kursrab}{\kursrabtrue}
\DeclareOption{remarks}{\remarkstrue}
\DeclareOption{hyperprint}{\hyperprinttrue}
\DeclareOption{greekit}{\greekittrue}
\DeclareOption{hyperref}{\hyperreftrue}
\DeclareOption{nodotheadline}{\nodotheadlinetrue}


%% Метки на полях - чтобы было удобно замечания писать!

\marginparwidth 37mm
\marginparsep 2mm

% \oddsidemargin=-0.5cm
% \textwidth=17cm
% \topmargin=-1.5cm%0pt
% \textheight=25cm

 \oddsidemargin=0.46cm
 \textwidth=16.0cm
 \topmargin=-0.54cm%0pt
 \textheight=24cm

\newcommand{\rem}[1]{%
\ifremarks%
 \renewcommand{\baselinestretch}{1.0}
 \selectfont
 \hoffset=-25mm
 %\reversemarginpar
 \marginpar{\small\raggedright\sffamily#1}%
\else%
\fi%
}


\ProcessOptions


\RequirePackage{amssymb}
\RequirePackage{amsmath}
\RequirePackage{cite}
\RequirePackage{caption}
\RequirePackage{cmap}
%\RequirePackage{unicode}
\RequirePackage{bm}
\RequirePackage{graphicx}
\RequirePackage{indentfirst}
\RequirePackage{ifpdf}
\RequirePackage{soulutf8}

\ifpdf%
  %{
   \ifhyperref
   \usepackage{hyperref}%
   \usepackage[numbered]{bookmark}%
   \fi%
  %}
\fi


%%\ifgreekit
%%\def\omicron{o}
%%\else
%%%Прямые греческие
%% \let\Alpha\relax
%% \let\Beta\relax
%% \let\Gamma\relax
%% \let\Delta\relax
%% \let\Epsilon\relax
%% \let\Zeta\relax
%% \let\Eta\relax
%% \let\Theta\relax
%% \let\Iota\relax
%% \let\Kappa\relax
%% \let\Lambda\relax
%% \let\Mu\relax
%% \let\Nu\relax
%% \let\Xi\relax
%% \let\Omicron\relax
%% \let\Pi\relax
%% \let\Rho\relax
%% \let\Sigma\relax
%% \let\Tau\relax
%% \let\Upsilon\relax
%% \let\Phi\relax
%% \let\Chi\relax
%% \let\Psi\relax
%% \let\Omega\relax
%%
%% \DeclareFontFamily{T1}{umathgr}{}
%% \DeclareFontShape{T1}{umathgr}{m}{n}{%
%%  <-> omsegr }{}
%% \DeclareFontShape{T1}{umathgr}{b}{n}{%
%%  <-> omsegrb }{}
%% \DeclareFontShape{T1}{umathgr}{m}{it}{%
%%  <-> omsegri }{}
%% \DeclareSymbolFont{um@thgr}{T1}{umathgr}{m}{n}
%% \SetSymbolFont{um@thgr}{bold}{T1}{umathgr}{b}{n}
%% \DeclareMathAlphabet{\mathgbf}{T1}{umathgr}{b}{n}
%% \DeclareMathAlphabet{\mathgit}{T1}{umathgr}{m}{it}
%%
%% \DeclareMathSymbol{\alpha}{\mathalpha}{um@thgr}{"61}
%% \DeclareMathSymbol{\Alpha}{\mathalpha}{um@thgr}{"40}
%%
%% \DeclareMathSymbol{\beta}{\mathalpha}{um@thgr}{"62}
%% \DeclareMathSymbol{\varbeta}{\mathalpha}{um@thgr}{"31}
%% \DeclareMathSymbol{\Beta}{\mathalpha}{um@thgr}{"41}
%%
%% \DeclareMathSymbol{\gamma}{\mathalpha}{um@thgr}{"67}
%% \DeclareMathSymbol{\Gamma}{\mathalpha}{um@thgr}{"47}
%%
%% \DeclareMathSymbol{\delta}{\mathalpha}{um@thgr}{"64}
%% \DeclareMathSymbol{\Delta}{\mathalpha}{um@thgr}{"44}
%%
%% \DeclareMathSymbol{\varepsilon}{\mathalpha}{um@thgr}{"65}
%% \DeclareMathSymbol{\Epsilon}{\mathalpha}{um@thgr}{"45}
%%
%% \DeclareMathSymbol{\zeta}{\mathalpha}{um@thgr}{"7A}
%% \DeclareMathSymbol{\Zeta}{\mathalpha}{um@thgr}{"5A}
%%
%% \DeclareMathSymbol{\eta}{\mathalpha}{um@thgr}{"68}
%% \DeclareMathSymbol{\Eta}{\mathalpha}{um@thgr}{"48}
%%
%% \DeclareMathSymbol{\theta}{\mathalpha}{um@thgr}{"6A}
%% \DeclareMathSymbol{\vartheta}{\mathalpha}{um@thgr}{"32}
%% \DeclareMathSymbol{\Theta}{\mathalpha}{um@thgr}{"4A}
%%
%% \DeclareMathSymbol{\iota}{\mathalpha}{um@thgr}{"69}
%% \DeclareMathSymbol{\Iota}{\mathalpha}{um@thgr}{"49}
%%
%% \DeclareMathSymbol{\kappa}{\mathalpha}{um@thgr}{"35}
%% \DeclareMathSymbol{\varkappa}{\mathalpha}{um@thgr}{"6B}
%% \DeclareMathSymbol{\Kappa}{\mathalpha}{um@thgr}{"4B}
%%
%% \DeclareMathSymbol{\lambda}{\mathalpha}{um@thgr}{"6C}
%% \DeclareMathSymbol{\Lambda}{\mathalpha}{um@thgr}{"4C}
%%
%% \DeclareMathSymbol{\mu}{\mathalpha}{um@thgr}{"6D}
%% \DeclareMathSymbol{\Mu}{\mathalpha}{um@thgr}{"4D}
%%
%% \DeclareMathSymbol{\nu}{\mathalpha}{um@thgr}{"6E}
%% \DeclareMathSymbol{\Nu}{\mathalpha}{um@thgr}{"4E}
%%
%% \DeclareMathSymbol{\xi}{\mathalpha}{um@thgr}{"78}
%% \DeclareMathSymbol{\Xi}{\mathalpha}{um@thgr}{"51}
%%
%% \DeclareMathSymbol{\omicron}{\mathalpha}{um@thgr}{"6F}
%% \DeclareMathSymbol{\Omicron}{\mathalpha}{um@thgr}{"4F}
%%
%% \DeclareMathSymbol{\pi}{\mathalpha}{um@thgr}{"70}
%% \DeclareMathSymbol{\varpi}{\mathalpha}{um@thgr}{"37}
%% \DeclareMathSymbol{\Pi}{\mathalpha}{um@thgr}{"50}
%%
%% \DeclareMathSymbol{\rho}{\mathalpha}{um@thgr}{"72}
%% \DeclareMathSymbol{\varrho}{\mathalpha}{um@thgr}{"34}
%% \DeclareMathSymbol{\Rho}{\mathalpha}{um@thgr}{"52}
%%
%% \DeclareMathSymbol{\sigma}{\mathalpha}{um@thgr}{"73}
%% \DeclareMathSymbol{\varsigma}{\mathalpha}{um@thgr}{"62}
%% \DeclareMathSymbol{\Sigma}{\mathalpha}{um@thgr}{"53}
%%
%% \DeclareMathSymbol{\tau}{\mathalpha}{um@thgr}{"74}
%% \DeclareMathSymbol{\Tau}{\mathalpha}{um@thgr}{"54}
%%
%% \DeclareMathSymbol{\upsilon}{\mathalpha}{um@thgr}{"75}
%% \DeclareMathSymbol{\Upsilon}{\mathalpha}{um@thgr}{"55}
%% \DeclareMathSymbol{\Varupsilon}{\mathalpha}{um@thgr}{"30}
%%
%% \DeclareMathSymbol{\phi}{\mathalpha}{um@thgr}{"33}
%% \DeclareMathSymbol{\varphi}{\mathalpha}{um@thgr}{"66}
%% \DeclareMathSymbol{\Phi}{\mathalpha}{um@thgr}{"46}
%%
%% \DeclareMathSymbol{\chi}{\mathalpha}{um@thgr}{"71}
%% \DeclareMathSymbol{\Chi}{\mathalpha}{um@thgr}{"51}
%%
%% \DeclareMathSymbol{\psi}{\mathalpha}{um@thgr}{"79}
%% \DeclareMathSymbol{\Psi}{\mathalpha}{um@thgr}{"59}
%%
%% \DeclareMathSymbol{\omega}{\mathalpha}{um@thgr}{"77}
%% \DeclareMathSymbol{\Omega}{\mathalpha}{um@thgr}{"57}
%%
%%\fi

%\RequirePackage{geometry}






\pagestyle{headings}

\makeatletter

%Русские переносы - 2 знака от конца
\righthyphenmin=2
\frenchspacing

%Перенос знака математической операции на новую строчку
\newcommand\brop[1]{#1\discretionary{}{\hbox{$#1$}}{}}

%Красивые знаки "больше либо равно" и "меньше либо равно"
\renewcommand{\ge}{\geqslant}
\renewcommand{\le}{\leqslant}

%Пустое множество
\def\emptyset{\varnothing}

%Интерлиньяж (междустрочный интервал)
\renewcommand{\baselinestretch}{1.3}

%%Изменение отбивок в перечнях на нормальные
\renewcommand{\@listI}{%
\leftmargin=32pt \rightmargin=0pt \labelsep=5pt \labelwidth=10pt
\itemindent=0pt \listparindent=0pt%%
\topsep=1pt plus 1pt minus 1 pt
\partopsep=2pt plus 1pt minus 1pt
\parsep=0pt plus 1pt
\itemsep=\parsep}

\renewcommand{\@listii}{%
\leftmargin=25pt \rightmargin=0pt \labelsep=5pt \labelwidth=20pt
\itemindent=0pt \listparindent=0pt%%
\topsep=1pt plus 1pt minus 1 pt
\partopsep=2pt plus 1pt minus 1pt
\parsep=0pt plus 1pt
\itemsep=\parsep}

%% Команда центрирования для многострочных заголовков
\def\raggedcenter@{\leftskip\z@ plus.4\hsize \rightskip\leftskip
 \parfillskip\z@ \parindent\z@ \spaceskip.5em \xspaceskip.7em\relax
 \pretolerance9999\tolerance9999 \exhyphenpenalty\@M\relax
 \hyphenpenalty\@M \let\\\linebreak}

%% Точка после номеров рисунков и таблиц
\captionsetup{labelsep=period}



%\newcommand*\stretchwith[2]{%
%  \let\@tempa\relax
%  \def\@tempb{#1}%
%  \@tfor\@tempc:=#2\do {\@tempa \@tempc \let\@tempa\@tempb}%
%}



\ifnodotheadline
 \def\@@dothead{\ }
\else
 \def\@@dothead{.\ }
\fi


%Формируем колонтитулы: в колонтитулы название раздела (либо то, что стоит в квадратных скобках)
\def\sectionmark#1{\markboth%
{\ifnum\c@secnumdepth>\z@\conthdsection\thesection\@@dothead\ \fi#1}%
{\ifnum\c@secnumdepth>\z@\conthdsection\thesection\@@dothead\ \fi#1}%
}


%%-----------------------------------------------------------
%% Аналог команды @startsection, но с точками после номеров.
%%-----------------------------------------------------------
\def\@startctsection#1#2#3#4#5#6{\if@noskipsec \leavevmode \fi
   \par \@tempskipa #4\relax
   \@afterindenttrue
   \ifdim \@tempskipa <\z@ \@tempskipa -\@tempskipa \@afterindentfalse\fi
   \if@nobreak \everypar{}\else
     \addpenalty{\@secpenalty}\addvspace{\@tempskipa}\fi
     \@ifstar
     %%%{\@sctsect{#3}{#4}{#5}{#6}}                                         %%% ---  И.К. 29.09.2009
     {\@sctsect{0mm}{#4}{#5}{#6}}                                           %%% ---  И.К. 29.09.2009
     {\@threarg{\@ctsect{#1}{#2}{#3}{#4}{#5}{#6}}}}

\long\def\@threarg#1{\@ifnextchar [{\@nthrearg{#1}}{\@xthrearg{#1}}}
\long\def\@xthrearg#1#2{#1[{#2}][{#2}]{#2}}
\long\def\@nthrearg#1[#2]{\@ifnextchar [{#1[#2]}{#1[#2][#2]}}

%\long\def\@dblarg#1{\@ifnextchar[{#1}{\@xdblarg{#1}}}
%\long\def\@xdblarg#1#2{#1[{#2}]{#2}}

\def\@ctsect#1#2#3#4#5#6[#7][#8]#9{\ifnum #2>\c@secnumdepth
     \let\@svsec\@empty\else
     \refscstepcounter{#1}\def\@svsec{\@makectsecthead{#1}\hskip #3}\fi
     \ifdefined\href
        \Hy@GlobalStepCount\Hy@linkcounter                                         % И.К.
        \xdef\@currentHref{section*.\the\Hy@linkcounter}%                          % И.К.
        \Hy@raisedlink{\hyper@anchorstart{\@currentHref}\hyper@anchorend}%         % И.К.
       %
     \else
        %\message{no package hyperref}%
     \fi

     \@tempskipa #5\relax
      \ifdim \@tempskipa>\z@
        \begingroup
        #6\relax
          \ifleftheadline{}
           \@hangfrom{\@svsec\hskip #3\relax                                 %%%    --- И.К. 29.09.2009
                    }{\interlinepenalty \@M #9\par}%                         %%%    --- И.К. 29.09.2009
          \else {}
           \raggedcenter@\interlinepenalty\@M #6\relax\@svsec #9\par         %%%    --- И.К. 29.09.2009
          \fi

        \endgroup
       \csname #1mark\endcsname{#8}\addcontentsline
         {toc}{#1}{\ifnum #2>\c@secnumdepth \else
                      \protect\numberline{%
                      \ifpdf\ifhyperref%
                      \texorpdfstring%
                      {\hbox to 0pt{\hss\csname conthd#1\endcsname}}
                      {\hss \csname conthd#1\endcsname}%
                      \else%
                      {\hbox to 0pt{\hss\csname conthd#1\endcsname}}
                      \fi\fi%
                      %
                      \ifboldsect{%
                      \ifnum #2=1 {%
                      \textbf{\csname the#1\endcsname\@@dothead}%
                      }%
                      \else%
                      {\csname the#1\endcsname\@@dothead}\fi%
                      }%
                      \else%
                      {{\csname the#1\endcsname\@@dothead}}%
                      \fi%
                      %
                      }\fi%
                    \ifboldsect{%
                    \ifnum #2=1{%
                    %\textbf{#7}%                                    %%%    --- И.К. 23.05.2011
                    %%%%%\textbf{\ignorespaces#7}%                   %%%    --- И.К. 23.05.2011
                    \textbf{\ignorespaces\ifx&#7&{#9}\else{#7}\fi}%  %%%    --- И.К. 29.09.2011
                    }%
                    \else%
                    %#7%                                        %%%    --- И.К. 23.05.2011
                    %%%%%\ignorespaces#7%                       %%%    --- И.К. 23.05.2011
                    \ignorespaces\ifx&#7&{#9}\else{#7}\fi%      %%%    --- И.К. 29.09.2011
                    \fi%
                    }%
                    \else%
                    %#7%                                        %%%    --- И.К. 23.05.2011
                    %%%%%\ignorespaces#7%                       %%%    --- И.К. 23.05.2011
                    \ignorespaces\ifx&#7&{#9}\else{#7}\fi%      %%%    --- И.К. 29.09.2011
                    \fi%
                    }\else
        \def\@svsechd{#6\hskip #3\relax  %% \relax added 2 May 90
                   \@svsec #9\csname #1mark\endcsname
                      {#8}\addcontentsline
                           {toc}{#1}{\ifnum #2>\c@secnumdepth \else
                             \protect\numberline{\csname the#1\endcsname\@@dothead}\fi
                       %#7}}\fi                                 %%%    --- И.К. 23.05.2011
                       \ignorespaces#7}}\fi                     %%%    --- И.К. 23.05.2011
     \@xsect{#5}}



\def\@xsect#1{\@tempskipa #1\relax
      \ifdim \@tempskipa>\z@
       \par \nobreak
       \vskip \@tempskipa
       \@afterheading
    \else \global\@nobreakfalse \global\@noskipsectrue
       \everypar{\if@noskipsec \global\@noskipsecfalse
                   \clubpenalty\@M \hskip -\parindent
                   \begingroup \@svsechd \endgroup \unskip
                   \hskip -#1\relax  % relax added 14 Jan 91
                  \else \clubpenalty \@clubpenalty
                    \everypar{}\fi}\fi\ignorespaces}

\def\@sctsect#1#2#3#4#5{\@tempskipa #3\relax
   \ifdim \@tempskipa>\z@
     {\begingroup
   \ifleftheadline{}
     \interlinepenalty\@M\noindent#4\hskip#1\relax#5\par
   \else
     \interlinepenalty\@M \raggedcenter@ #4\hskip #1\relax #5\par     %%%    --- И.К. 29.09.2009
   \fi
     \endgroup}
   \else
     \begingroup
      \ifleftheadline{}
       \def\@svsechd{\noindent#4\hskip #1\relax #5}            %%%    --- И.К. 29.09.2009
      \else
       \def\@svsechd{\raggedcenter@ #4\hskip #1\relax #5}            %%%    --- И.К. 29.09.2009
      \fi
     \endgroup
     \fi
    \@xsect{#3}}

%%%\def\@makectsecthead#1{\csname head#1 \endcsname \csname the#1 \endcsname.}
\def\@makectsecthead#1{\csname head#1\endcsname\upshape{\csname the#1\endcsname\@@dothead\ignorespaces}} %%% И.К. 23.05.2011

\newcommand{\sectionfont}{\Large\bf}
\def\section{   \@ifnextchar -{\@@nsection{\secsepbefore}{\secsepafter}
                \reset@font\sectionfont}{\@startctsection
                {section}{1}{\@secdelim}
                {\secsepbefore}{\secsepafter}{\reset@font\sectionfont}}}
\def\secsepbefore{3.5ex plus 1ex minus .2ex}
\def\secsepafter{2.3ex plus.2ex}
\def\@@nsection#1#2#3-{\if@noskipsec \leavevmode \fi
   \par \@tempskipa #1\relax
   \@afterindenttrue
   \ifdim \@tempskipa <\z@ \@tempskipa -\@tempskipa \@afterindentfalse\fi
   \if@nobreak \everypar{}\else
     \addpenalty{\@secpenalty}\addvspace\@tempskipa\fi
    \@threarg{\@nsection{#1}{#2}{#3}}}
\def\@nsection#1#2#3[#4][#5]#6{\relax
     \@tempskipa #2\relax
        \begingroup
        \ifdefined\href
          \Hy@GlobalStepCount\Hy@linkcounter                                         % И.К.
          \xdef\@currentHref{section*.\the\Hy@linkcounter}%                          % И.К.
          \Hy@raisedlink{\hyper@anchorstart{\@currentHref}\hyper@anchorend}%         % И.К.
          %
        \else
          %\message{no package hyperref}%
         \fi

        \ifleftheadline{}
         #3\noindent#6\par                                %%% И.К. - 29.09.2009
        \else
         \raggedcenter@ #3 #6\par                         %%% И.К. - 29.09.2009
        \fi

        \endgroup
        \markright{#5}\relax

        \ifboldsect{
        \addcontentsline
         %%%{toc}{section}{\protect\numberline{}#5}%         %%% И.К. - 29.09.2009
         %%%{toc}{section}{\textbf{#5}}%                     %%% И.К. - 29.09.2009
         %%%%%{toc}{section}{\textbf{\ignorespaces#5}}%      %%% И.К. - 23.05.2011
         {toc}{section}{\textbf{\ignorespaces\ifx&#4&{#6}\else{#4}\fi}}%      %%% И.К. - 29.09.2011
         \ifpdf\ifhyperref%
          \phantomsection%
         \fi\fi
         %} \else { \addcontentsline{toc}{section}{#5} } \fi                  %%% И.К. - 23.05.2011
         %%%%%} \else { \addcontentsline{toc}{section}{\ignorespaces#5} } \fi %%% И.К. - 23.05.2011
         } \else { \addcontentsline{toc}{section}{\ignorespaces\ifx&#4&{#6}\else{#4}\fi} } \fi %%% И.К. - 29.09.2011
     \@xsect{#2}}

\def\secdelim#1{\def\@secdelimamount{#1}\gdef\@secdelim{\@secdelimamount}}
\def\refscstepcounter#1{\stepcounter{#1}\let\@tempa\protect
\def\protect{\noexpand\protect\noexpand}%
\edef\@currentlabel{\csname p@#1\endcsname
                    \csname conthd#1\endcsname
                    \csname the#1\endcsname}%
\let\protect\@tempa}

\gdef\conthdchapter{}\gdef\appendixname{Приложение}
\gdef\headsection{}\gdef\conthdsection{}
\secdelim{0.1em}
%%%\def\thesection{\thechapter.\arabic{section}}                  %%%% И.К. - 29.09.2009
\def\thesection{\arabic{section}}                                 %%%% И.К. - 29.09.2009
\def\contentsname{Оглавление}
%%%\def\subsection{\@startctsection {subsection}{2}{0.3em}%     %%%% И.К. - 29.09.2009
\def\subsection{\@startctsection {subsection}{2}{\@secdelim}%     %%%% И.К. - 29.09.2009
        {3ex plus 1ex minus .2ex}{2ex plus 0.3ex minus 0.1ex}{\bf}}
\def\subsubsection{\@startctsection {subsubsection}{3}{\@secdelim}%                %%%% И.К. - 07.12.2013
        {3ex plus 1ex minus .2ex}{2ex plus 0.3ex minus 0.2ex}{\bf}}     %%%% И.К. - 29.09.2009
\def\headsubsection{}\def\conthdsubsection{}
\def\headsubsubsection{}\def\conthdsubsubsection{}				   %%%% И.К. - 07.12.2013	
%%%\def\l@section{\@dottedtocline{1}{1.7em}{2.3em}}                     %%%% И.К. - 29.09.2009
\def\l@section{\vspace*{3mm}\@dottedtocline{1}{0em}{1.6em}}
\def\l@subsection{\vspace*{2mm}\@dottedtocline{2}{1.6em}{2.2em}}
\def\l@subsubsection{\vspace*{1mm}\@dottedtocline{2}{3.8em}{2.8em}}
\newcommand\initcounter[1]{\stepcounter{#1}\addtocounter{#1}{-1}}





%%%%%%% Нумерация формул - по разделам
\ifeqwhole \else \renewcommand\theequation{\arabic{section}.\arabic{equation}} \@addtoreset{equation}{section}\fi


%%%%%%% Нумерация рисунков - по разделам
\iffigwhole \else \renewcommand\thefigure{\arabic{section}.\arabic{figure}} \@addtoreset{figure}{section}\fi


%%%%%%% Сноски - добавляем пробел
\renewcommand{\@makefntext}[1]{\renewcommand{\baselinestretch}{1.0}\selectfont\hangindent=2.3em\hangafter=1\parindent=1em\noindent\hbox to 2.3em{\hss\@makefnmark\rule{0.5em}{0pt}}#1}
%\@addtoreset{footnote}{page}
\iffootwhole
\else
\@addtoreset{footnote}{page}
\fi


\ifoglav{\addto\captionsrussian{\def\contentsname{Оглавление}}}\fi

\ifleftheadline{}\else{ \def\tableofcontents{
\markboth{\contentsname}{\contentsname}
\thispagestyle{plain}
\section*{\protect{\hbox to \textwidth{\hfill \contentsname \hfill}}}\@starttoc{toc} }}\fi


\renewcommand{\@biblabel}[1]{#1.\hfill}

%%%\addto\captionsrussian{\def\refname{Список литературы}}

\ifspisok{\addto\captionsrussian{\def\refname{Список использованных источников}}}\else{\addto\captionsrussian{\def\refname{Литература}}}\fi


%%%        Оформление списка литературы и цитирования источников

\renewenvironment{thebibliography}[1]
     {

      \ifleftheadline{\section-{\refname}} \else{\section-[\refname]{\protect{\hbox to \textwidth{\hfill \refname \hfill}}}} \fi

      %\section-{\bibname}%    %   Замена \chapter* на chapter- для вставки в оглавление. АК
      \@mkboth{\refname}{\refname}%
      \list{\@biblabel{\@arabic\c@enumiv}}%
           {\settowidth\labelwidth{\@biblabel{#1}}%
            \leftmargin\labelwidth
            \advance\leftmargin\labelsep
            \@openbib@code
            \usecounter{enumiv}%
            \let\p@enumiv\@empty
            \renewcommand\theenumiv{\@arabic\c@enumiv}}%
      \sloppy
      \clubpenalty4000
      \@clubpenalty \clubpenalty
      \widowpenalty4000%
      \sfcode`\.\@m}
     {\def\@noitemerr
       {\@latex@warning{Empty `thebibliography' environment}}%
      \endlist}


%%% Колонтитулы





   \def\@oddhead{\vbox to \headheight{%
                \hbox to \textwidth{%
                   \hfil{\footnotesize\rightmark}\hfil
                     \llap{\rm\thepage}}%
                     \kern2pt\hrule\vss}}%


% \def\sectionmark#1{\markboth{}{\uppercase{\ifnum \c@secnumdepth >\z@
%   \conthdsection\thesection.\ \fi #1}}}}


%% Обложка

\def\group#1{\def\@group{#1}}
\def\supervisor#1{\def\@supervisor{#1}} 
\def\supervisorR#1{\def\@supervisorR{#1}} 

\renewcommand{\maketitle}{
\thispagestyle{empty}
{
\centering

\renewcommand{\baselinestretch}{1.0}
\selectfont

\vspace*{-10mm}

\begin{minipage}{30mm}
\ifpdf
 \includegraphics[scale=0.14]{emblema}
\fi
\end{minipage}
\begin{minipage}{127mm}
\begin{center}
\fontsize{9}{11pt}\selectfont\textbf{Министерство науки и высшего образования Российской Федерации}

\fontsize{9}{11pt}\selectfont\textbf{Федеральное государственное бюджетное образовательное учреждение}

\fontsize{9}{11pt}\selectfont\textbf{высшего образования}

\fontsize{9}{11pt}\selectfont\textbf{<<Московский государственный технический университет}

\fontsize{9}{11pt}\selectfont\textbf{имени Н.\,Э.~Баумана}

\fontsize{9}{11pt}\selectfont\textbf{(национальный исследовательский университет)>>}

\fontsize{9}{11pt}\selectfont\textbf{(МГТУ им.~Н.\,Э.~Баумана)}
\end{center}
\end{minipage}

\vspace*{2mm}

\hrule height .7mm



\vspace*{2pt}

\hrule

}


\vspace*{9mm}

\noindent ФАКУЛЬТЕТ \rlap{\rule[-1.2mm]{132mm}{0.2mm}}\hbox to 132mm{\hfill Фундаментальные науки \hfill}

\medskip

\noindent КАФЕДРА \rlap{\rule[-1.2mm]{137.5mm}{0.2mm}}\hbox to 137.5mm{\hfill Прикладная математика \hfill}

\vspace*{2cm}

\centerline{\fontsize{18}{20pt}\selectfont\textbf{РАСЧЕТНО-ПОЯСНИТЕЛЬНАЯ ЗАПИСКА}}

\vspace*{5mm}

\centerline{\fontsize{17}{20pt}\selectfont\textit{\textbf{К КУРСОВОЙ РАБОТЕ}}}

\vspace*{5mm}

\centerline{\fontsize{17}{20pt}\selectfont\textit{\textbf{НА ТЕМУ:}}}

\begin{center}
\fontsize{18}{24pt}\selectfont\textit{\textbf{\@title}}
\end{center}

\vfill

\noindent\begin{minipage}{55mm}
\small
\noindent Студент \rlap{\rule[-1.2mm]{39mm}{0.2mm}}\hbox to 39mm{\hfill \@group \hfill}

\noindent \phantom{Студент} \raisebox{2mm}{\hbox to 39mm{\hfill \footnotesize(Группа) \hfill}}

\vspace*{3.5mm}

\noindent{Руководитель курсовой работы}

\vspace*{10mm}

\noindent{Руководитель курсовой работы}

\vspace*{5mm}

\end{minipage}
\rule{20mm}{0pt}
\begin{minipage}{79mm}
\noindent\rlap{\rule[-1.2mm]{30mm}{0.2mm}}\hbox to 30mm{\hfill} \quad \rlap{\rule[-1.2mm]{47mm}{0.2mm}}{\rule{3mm}{0pt}\@author}

\noindent\raisebox{2mm}{\hbox to 30mm{\hfill \footnotesize(Подпись, дата) \hfill}} \quad  \raisebox{2mm}{\rule{5mm}{0pt}\footnotesize(И.\,О.~Фамилия)}

\vspace*{2mm}

\noindent\rlap{\rule[-1.2mm]{30mm}{0.2mm}}\hbox to 30mm{\hfill} \quad \rlap{\rule[-1.2mm]{47mm}{0.2mm}}{\rule{3mm}{0pt}\@supervisor}

\noindent\raisebox{2mm}{\hbox to 30mm{\hfill \footnotesize(Подпись, дата) \hfill}} \quad  \raisebox{2mm}{\rule{5mm}{0pt}\footnotesize(И.\,О.~Фамилия)}

\vspace*{2mm}

\noindent\rlap{\rule[-1.2mm]{30mm}{0.2mm}}\hbox to 30mm{\hfill} \quad \rlap{\rule[-1.2mm]{47mm}{0.2mm}}{\rule{3mm}{0pt}\@supervisorR}

\noindent\raisebox{2mm}{\hbox to 30mm{\hfill \footnotesize(Подпись, дата) \hfill}} \quad  \raisebox{2mm}{\rule{5mm}{0pt}\footnotesize(И.\,О.~Фамилия)}

\end{minipage}

\vfill

\centerline{\large\@date\ г.}

\pagebreak}


%% Старая обложка - автор Пузикова В.В.

\newcommand{\oldcover}[9]{
\thispagestyle{empty}
{
\centering


\raisebox{15mm}[1mm]{\textrm{М\,и\,н\,и\,с\,т\,е\,р\,с\,т\,в\,о\,\,\,\,о\,б\,р\,а\,з\,о\,в\,а\,н\,и\,я\,\,\,\,и\,\,\,\,н\,а\,у\,к\,и\,\,\,\,Р\,о\,с\,с\,и\,й\,с\,к\,о\,й\,\,\,\,Ф\,е\,д\,е\,р\,а\,ц\,и\,и}}

\kern-12mm
{
\fontsize{11}{14pt}\selectfont
\textbf{Федеральное государственное бюджетное образовательное учреждение\\высшего профессионального образования\\%
<<Московский\ государственный\ технический\ университет\ имени\ Н.\,Э.\,Баумана>>}\\[3mm]
\hrule height .7mm

}

\vspace*{2pt}

\hrule

\bigskip%
{\large\textit{Кафедра <<Прикладная математика>>}}\\%
\vspace*{\fill}%
\includegraphics[scale=0.4]{emblema}\\%
\vspace*{\fill}%
\Large\textit{
\ifkursrab\textbf{Курсовая работа}
\else\textbf{Курсовой проект}\fi
}\\[2mm]%
\ifx&#1&%
\else
 \textrm{по дисциплине <<#1>>}\\[2mm]
\fi%
\Large{\textbf{#2}}%
%\end{center}%

}
\vspace*{\fill}%
{
\centering

\bigskip

%\begin{center}%
\begin{tabular}{@{}p{3.0cm}@{}p{7.7cm}@{}p{0.3cm}@{}p{5.0cm}@{}}%
\textit{Выполнил}& {#3} && \textit{#4}\\%
\bigskip\\%
\parbox[t]{3.0cm}{\textit{Научный}\\[-1pt] \textit{руководитель}}& \parbox[t]{7.7cm}{{#5}} && \textit{#6}\\%
\bigskip\\%
\ifx&#7&%
\else
 \textit{Консультант}& \parbox[t]{7.7cm}{{#7}} && \textit{#8}\\
\fi%
\end{tabular}\\%
\vspace*{\fill}
Москва --- {#9}\\
%\hrulefill%\end{center}

}
\pagebreak}






%Настройки вида гиперссылок

\ifhyperprint{}
\ifdefined\href
\hypersetup{unicode=true, colorlinks=true, linkcolor=black, anchorcolor=black, citecolor=black, filecolor=black, menucolor=black, runcolor=black, urlcolor=black, bookmarks=true}
\fi
\else
\ifdefined\href
\hypersetup{unicode=true, colorlinks=true, bookmarks=true}
\fi
\fi


%%Полужирный наклонный шрифт в формулах
\newcommand\be[1]{\text{\mathversion{bold}${#1}$}}
\newcommand\bi[1]{\text{\mathversion{bold}${#1}$}}
\newcommand\mathbi[1]{\text{\mathversion{bold}${#1}$}}

\newcommand\mi[1]{\text{\mathversion{normal}${#1}$}}
\newcommand\mathmi[1]{\text{\mathversion{normal}${#1}$}}

%%Дробь среднего размера \sfrac
\def\sfrac#1#2{\hbox{\footnotesize$\displaystyle\frac{#1}{#2}$}}

%%        Текст вразрядку. Обращение \looser{<число>}{текст}
%%        <число> - величина клея в em'ах между буквами.
\def\looser#1#2{%
\let\@mysp\relax%
\let\@myspinner\relax%
\newlength{\@mysp}%
\newlength{\@myspinner}%
\setlength{\@mysp}{#1em}%
\ifdim\@mysp=0em%
\fi%
\setlength{\@myspinner}{2.0\@mysp}%
\addtolength{\@myspinner}{0.333em plus.03em minus.03em}%
\sodef\textso{}{\@mysp}{\@myspinner}{\@myspinner}%
\textso{#2}%
\resetso%
}


%\input letterspacing.tex
%\newcommand{\kerning}[2]{\mbox{\letterspace to #1\naturalwidth{#2}}}


%\input{letterspacing.tex}\newcommand{\tracking}[2]{\mbox{\letterspace to #1\naturalwidth{#2}}}

%\newcount\@oldcatcode
%\def\looser{\@oldcatcode=\the\catcode` \catcode` =13\relax\@looser}
%\def\@looser#1#2{\bgroup
%     \def\loos@space{\kern #1em}%
%     \def\loos@next{\loos@text}%
%     \def\loos@text ##1{##1\loos@space\loos@next}%
%     \loos@text #2{\let\loos@next=\relax\catcode` =\the\@oldcatcode}%
%     \egroup}

%\def\@looser#1#2{\def\loos@space{\kern #1em}%
%     \def\loos@text ##1{##1\loos@space\noexpand\loos@text}%
%     \loos@text #2{\kern-#1em\let\loos@text=\relax
%                  \catcode` =\the\@oldcatcode}\relax}

%\def\@looser#1#2{#2}

%% Работа с рисунками


\def\DefaultPicExt#1{\def\@DefaultPicext{#1}}
\def\DefaultPath#1{\def\@DefaultPath{#1}}
\DefaultPicExt{}
\DefaultPath{}


\let\@figcollect\relax
\def\demofig{\ifANSI
                \typeout{LўхЄ Ёшёєэъют эх тхфхЄё }%
             \else
                \typeout{Учет рисунков не ведется}
             \fi}
\def\figcollect{\newcounter{allfig}%
        \def\@figcollect{\addtocounter{allfig}1}%
            \ifANSI
              \def\demofig{\typeout{TёЄЁхўхэю Ёшёєэъют: \arabic{allfig}}}%
            \else
              \def\demofig{\typeout{Встречено рисунков: \arabic{allfig}}}%
            \fi}


\def\@ParseFilename#1.#2;{\def\@CurrFilename{#1}\def\@CurrFileext{#2}}
%\def\@AddFileext#1#2{\@ParseFilename#1.;\ifx\@CurrFileext\@empty
%                    \def\@CurrFile{#1.#2}\else\def\@CurrFile{#1}\fi}
\def\@AddFileext#1#2{\@ParseFilename#1.;\ifx\@CurrFileext
                    \def\@CurrFile{#1.#2}\else\def\@CurrFile{#1}\fi}
\def\@SetCurrfileSpec#1{\@AddFileext{#1}{\@DefaultPicext}%
                        \ifx\@DefaultPath\@empty
                            \def\@CurrFileSpec{\@CurrFile}\else
                            \def\@CurrFileSpec{\@DefaultPath/\@CurrFile}\fi}

%\def\@PicSpecial#1#2{\edef\@SpecBody{{em:graph \@CurrFileSpec,#1,#2}}%
%                \let\@CurrFileSpec\relax
%                \expandafter\special\@SpecBody\relax}%

\def\@PicSpecial#1#2{\edef\@SpecBody{\@CurrFileSpec}%
                \let\@CurrFileSpec\relax%
                \expandafter\includegraphics[width=#1,height=#2]{\@SpecBody}\relax}%


\def\Grdunit#1{\def\@GrdU{#1}}
\Grdunit{pt}
%%  #1 ширина, #2 высота, #3 глубина (входит в высоту), #4 имя файла
\def\@gGraph#1#2#3#4{%
            \@tempdima=#1\relax
            \@tempdima=\@Graphscalefactor\@tempdima\relax
            \edef\@Graphwidth{\the\@tempdima}%
            \@tempdima=#2\@tempdima=\@Graphscalefactor\@tempdima
            \edef\@Graphheight{\the\@tempdima}%
            \if@permscale\else\gdef\@Graphscalefactor{1}\fi
            \setbox\@tempboxa
            \vbox to 0pt{\hsize=\@Graphwidth\offinterlineskip
            \hbox to \@Graphwidth{\@SetCurrfileSpec{#4}%
                        \@PicSpecial{\@Graphwidth}{\@Graphheight}\hfil}%
            \vss}%
            \@tempdima=-#3\advance\@tempdima by \@Graphheight
            \ht\@tempboxa=\@tempdima\dp\@tempboxa=#3
            \box\@tempboxa}

\def\@addfigfile#1#2{}

\def\@Graphscalefactor{1}
\newif\if@permscale\@permscalefalse

\def\Graph#1#2#3#4{\@addfigfile{graph}{#4}%
                   \@gGraph{#1}{#2}{#3}{#4}}


\def\picbox#1#2#3#4{\@addfigfile{picbox}{#3}%
                    \mbox{\vbox{\@tempdima=#1\relax
                          \@tempdima=\@Graphscalefactor\@tempdima
                          \hsize=\@tempdima\offinterlineskip
                          \hbox{\@gGraph{#1}{#2}{\@Grdepth}{#3}}%
                          \ifx{#4}%
                          \else
                             \vskip\@picboxsep%
                             \centering\captionfont #4%
                          \fi}}}
\def\setGrdepth#1{\def\@Grdepth{#1\gdef\@Grdepth{\Grdepth}}\ignorespaces}
\def\@Grdepth{0pt}

\def\figbox#1#2#3#4{\@ifnextchar[
                      {\@figbox{#1}{#2}{#3}{#4}}%
                      {\@figbox{#1}{#2}{#3}{#4}[]}}
\def\@figbox#1#2#3#4[#5]{\@addfigfile{figbox}{#3}%
               \mbox{\vbox{\@tempdima=#1\relax
                     \@tempdima=\@Graphscalefactor\@tempdima
                     \hsize=\@tempdima
                     \edef\@saveinterlineskip{\lineskip\the\lineskip
                                             \lineskiplimit\the\lineskiplimit}%
                     \offinterlineskip
                     \hbox{\@gGraph{#1}{#2}{0pt}{#3}}%
                     \vskip\@captionsep
                     \figcap[#5]{#4}}}}
\let\@saveinterlineskip\relax


%%%     Подписи к рисункам

\def\figcap{\@ifnextchar[{\@figcap}{\@figcap[]}}
\def\@figcap[#1]#2{\vbox{\@figcollect\@saveinterlineskip
                    \def\@captype{figure}\relax\caption{#1}\label{#2}}}


\long\def\@makecaption#1#2{%
  \def\@tempa{#2}\def\@tempb{\ignorespaces}%
  \ifx\@tempa\@tempb
     \def\@tempb{{#1}}%
  \else
     \def\@tempb{{#1.} }%
  \fi
  \sbox\@tempboxa{\captionfont\@tempb#2}%
  \ifdim \wd\@tempboxa >\hsize
     {\captionfont\@tempb#2\par}%
  \else
    \global \@minipagefalse
    \hb@xt@\hsize{\hfil\box\@tempboxa\hfil}%
  \fi}

%\newcommand\captionfont{\fontsize{9pt}{10pt}\fontshape{n}\fontseries{m}\selectfont}
\def\captionsep#1{\def\@captionsep{#1}\ignorespaces}
\def\picboxsep#1{\def\@picboxsep{#1}\ignorespaces}
\captionsep{10\p@}
\let\@picboxsep\@captionsep

\def\dfb#1{\ifvmode\let\@tempa=\relax\else\let\@tempa=\vadjust\fi\@tempa{\vbox to 0pt{#1\vss}}}

\def\figphantom#1#2#3#4#5#6{\@ifnextchar[
                      {\@figphantom{#1}{#2}{#3}{#4}{#5}{#6}}%
                      {\@figphantom{#1}{#2}{#3}{#4}{#5}{#6}[]}}
\def\@figphantom#1#2#3#4#5#6[#7]{\dfb{\kern#1\hbox{\kern#2\figbox{#3}{#4}{#5}{#6}[#7]}}}


\def\figph#1#2#3{\@ifnextchar[{\@figph{#1}{#2}{#3}}{\@@figph{#1}{#2}{#3}}}
\def\@@figph#1#2#3{\dfb{\kern#1\hbox{\kern#2{#3}}}}
\def\@figph#1#2#3[#4]{\dfb{\kern#1\parbox[t][1pt]{#4}{\kern#2{#3}}}}


\def\wbreak{{\parfillskip=0pt\par\noindent}\hangafter=0}


\makeatother


\endinput
